% Generated by `taskgen course examples` from ecc.ecdh (seed 7).
% Do not edit: change examples.yaml and regenerate.
\begin{example}[ECDH на маленькой кривой]\label{ex:ecdlp-ecdh}
Алиса и Боб используют ECDH на кривой $E\colon y^2 = x^3 + 34x + 2$ над $\mathbb{F}_{41}$ с базовой точкой $P = (7, 38)$. Секретные ключи: $a = 5$ у Алисы, $b = 14$ у Боба. Вычислить открытые ключи $A = aP$, $B = bP$ и общий ключ $K$ двумя способами ($aB$ и $bA$) методом «удвоение"--~сложение». Что видит Ева?

\par\noindent\textit{Решение.} Двоичная запись множителя читается слева направо: $R \gets 2R$, при единичном бите $R \gets R + X$. Формулы: $\lambda = \frac{3x_1^2 + a}{2y_1}$ (удвоение), $\lambda = \frac{y_2 - y_1}{x_2 - x_1}$ (сложение), $x_3 = \lambda^2 - x_1 - x_2$, $y_3 = \lambda(x_1 - x_3) - y_1$; если $x_1 = x_2$, $y_1 = -y_2$, то сумма равна $\mathcal{O}$; $\mathcal{O}$ "--- нейтральный элемент. Всё по модулю $41$, $a_E = 34$.\par\noindent\textbf{Алиса}, $A = aP$. $5 = 101_2$, $P = (7, 38)$:\[\begin{array}{l|r|r|c}\text{операция} & k & \lambda & R = kP \\ \hline 2R & 2 & 4 & (2, 23) \\ 2R & 4 & 1 & (38, 23) \\ R + P & 5 & 22 & (29, 11) \\ \end{array}\]\par\noindent\textbf{Боб}, $B = bP$. $14 = 1110_2$, $P = (7, 38)$:\[\begin{array}{l|r|r|c}\text{операция} & k & \lambda & R = kP \\ \hline 2R & 2 & 4 & (2, 23) \\ R + P & 3 & 3 & (0, 24) \\ 2R & 6 & 40 & (1, 18) \\ R + P & 7 & 17 & (35, 19) \\ 2R & 14 & 21 & (2, 18) \\ \end{array}\]\par\noindent\textbf{Алиса}, $K = aB$. $5 = 101_2$, $B = (2, 18)$:\[\begin{array}{l|r|r|c}\text{операция} & k & \lambda & R = kB \\ \hline 2R & 2 & 40 & (38, 18) \\ 2R & 4 & 37 & (22, 0) \\ R + B & 5 & 36 & (1, 18) \\ \end{array}\]\par\noindent\textbf{Боб}, $K = bA$. $14 = 1110_2$, $A = (29, 11)$:\[\begin{array}{l|r|r|c}\text{операция} & k & \lambda & R = kA \\ \hline 2R & 2 & 10 & (1, 23) \\ R + A & 3 & 23 & (7, 3) \\ 2R & 6 & 37 & (2, 18) \\ R + A & 7 & 21 & (0, 24) \\ 2R & 14 & 40 & (1, 18) \\ \end{array}\]\par Результаты совпадают: $aB = bA = abP = (1, 18)$ (на практике ключом служит $x$-координата $1$). Порядок $P$ равен $16$. Ева видит $E$, $P$, $A = (29, 11)$, $B = (2, 18)$; чтобы найти $K$, ей достаточно решить ECDLP $A = xP$.\[A = (29, 11),\quad B = (2, 18),\quad K = (1, 18).\]
\end{example}
